\chapter{2021年全国乙卷(理)}

\section{单选题}

\begin{problem}
设 \(2(z+\bar z)+3(z-\bar z)=4+6\i\),则 \(z=\)(\quad)

\choices
    {\(1-2\i\)}
    {\(1+2\i\)}
    {\(1+\i\)}
    {\(1-\i\)}

\end{problem}

\begin{problem}
已知集合 \(S=\{s\mid s=2n+1,n\in\Z\}\),\(T=\{t\mid t=4n+1,n\in\Z\}\),则 \(S\cap T=\)(\quad)

\choices
    {\(\varnothing\)}
    {\(S\)}
    {\(T\)}
    {\(\Z\)}

\end{problem}

\begin{problem}
已知命题 \(p:\exists x\in \R\),\(\sin x<1\);命题 \(q:\forall x\in \R\),\(\e^{|x|}\ge 1\),则下列命题中为真命题的是(\quad)

\choices
    {\(p\wedge q\)}
    {\(\neg p\wedge q\)}
    {\(p\wedge \neg q\)}
    {\(\neg\left( p\vee q \right)\)}

\end{problem}

\begin{problem}
设函数 \(f(x)=\frac{1-x}{1+x}\),则下列函数中为奇函数的是(\quad)

\choices
    {\(f(x-1)-1\)}
    {\(f(x-1)+1\)}
    {\(f(x+1)-1\)}
    {\(f(x+1)+1\)}

\end{problem}

\begin{problem}
在正方体 \(ABCD-A_1B_1C_1D_1\) 中,\(P\) 为 \(B_1D_1\) 的中点,则直线 \(PB\) 与 \(AD_1\) 所成的角为(\quad)

\bitmapfigure[max width=6.1cm]{img/2021/national_paper_B_science/q5_fig1.png}

\choices
    {\(\frac{\pi}{2}\)}
    {\(\frac{\pi}{3}\)}
    {\(\frac{\pi}{4}\)}
    {\(\frac{\pi}{6}\)}

\end{problem}

\begin{problem}
将 \(5\) 名北京冬奥会志愿者分配到花样滑冰,短道速滑,冰球和冰壶 \(4\) 个项目进行培训,每名志愿者只分配到 \(1\) 个项目,每个项目至少分配 \(1\) 名志愿者,则不同的分配方案共有(\quad)

\choices
    {\(60\) 种}
    {\(120\) 种}
    {\(240\) 种}
    {\(480\) 种}

\end{problem}

\begin{problem}
把函数 \(y=f(x)\) 图像上所有点的横坐标缩短到原来的 \(\frac{1}{2}\) 倍,纵坐标不变,再把所得曲线向右平移 \(\frac{\pi}{3}\) 个单位长度,得到函数 \(y=\sin\left( x-\frac{\pi}{4} \right)\) 的图像,则 \(f(x)=\)(\quad)

\choices
    {\(\sin\left( -\frac{x}{2}-\frac{7\pi}{12} \right)\)}
    {\(\sin\left( \frac{x}{2}+\frac{\pi}{12} \right)\)}
    {\(\sin\left( 2x-\frac{7\pi}{12} \right)\)}
    {\(\sin\left( 2x+\frac{\pi}{12} \right)\)}

\end{problem}

\begin{problem}
在区间 \((0,1)\) 与 \((1,2)\) 中各随机取 \(1\) 个数,则两数之和大于 \(\frac{7}{4}\) 的概率为(\quad)

\choices
    {\(\frac{7}{9}\)}
    {\(\frac{23}{32}\)}
    {\(\frac{9}{32}\)}
    {\(\frac{2}{9}\)}

\end{problem}

\begin{problem}
魏晋时期刘徽撰写的《海岛算经》是关于测量的数学著作. 其中第一题是测量海岛的高.

如图,点 \(E\),\(H\),\(G\) 在水平线 \(AC\) 上,\(DE\) 和 \(FG\) 是两个垂直于水平面且等高的测量标杆的高度,称为``表高'',\(EG\) 称为``表距'',\(GC\) 和 \(EH\) 都称为``表目距''. \(GC\) 与 \(EH\) 的差称为``表目距的差'',则海岛的高 \(AB=\)(\quad)

\FigureLayoutDeclare{img/2021/national_paper_B_science/q9_fig1.png}{5.2cm}{13bp 17bp 37bp 11bp}
\bitmapfigure[float=inline,max width=5.2cm]{img/2021/national_paper_B_science/q9_fig1.png}

\choices
    {\(\frac{\text{表高}\times \text{表距}}{\text{表目距的差}}+\text{表高}\)}
    {\(\frac{\text{表高}\times \text{表距}}{\text{表目距的差}}-\text{表高}\)}
    {\(\frac{\text{表高}\times \text{表距}}{\text{表目距的差}}+\text{表距}\)}
    {\(\frac{\text{表高}\times \text{表距}}{\text{表目距的差}}-\text{表距}\)}

\end{problem}

\begin{problem}
设 \(a\ne 0\),若 \(x=a\) 为函数 \(f(x)=a(x-a)^2(x-b)\) 的极大值点,则

\choices
    {\(a<b\)}
    {\(a>b\)}
    {\(ab<a^2\)}
    {\(ab>a^2\)}

\end{problem}

\begin{problem}
设 \(B\) 是椭圆 \(C:\frac{x^2}{a^2}+\frac{y^2}{b^2}=1\)(\(a>b>0\))的上顶点,若 \(C\) 上的任意一点 \(P\) 都满足 \(|PB|\le 2b\),则 \(C\) 的离心率的取值范围是(\quad)

\choices
    {\(\left[ \frac{\sqrt2}{2},1 \right)\)}
    {\(\left[ \frac{1}{2},1 \right)\)}
    {\(\left( 0,\frac{\sqrt2}{2} \right]\)}
    {\(\left( 0,\frac{1}{2} \right]\)}

\end{problem}

\begin{problem}
设 \(a=2\ln 1.01\),\(b=\ln 1.02\),\(c=\sqrt{1.04}-1\),则(\quad)

\choices
    {\(a<b<c\)}
    {\(b<c<a\)}
    {\(b<a<c\)}
    {\(c<a<b\)}

\end{problem}

\section{填空题}

\begin{problem}
已知双曲线 \(C:\frac{x^2}{m}-y^2=1\)(\(m>0\))的一条渐近线为 \(\sqrt3 x+my=0\),则 \(C\) 的焦距为\fillinblank{}.
\end{problem}

\begin{problem}
已知向量 \(\bs a=(1,3)\),\(\bs b=(3,4)\),若 \((\bs a-\lambda \bs b)\perp \bs b\),则 \(\lambda=\fillinblank{}\).
\end{problem}

\begin{problem}
记 \(\triangle ABC\) 的内角 \(A\),\(B\),\(C\) 的对边分别为 \(a\),\(b\),\(c\),面积为 \(\sqrt3\),\(B=60^\circ\),\(a^2+c^2=3ac\),则 \(b=\fillinblank{}\).
\end{problem}

\begin{problem}
以图 \(\circled{1}\) 为正视图,在图 \(\circled{2}\),\(\circled{3}\),\(\circled{4}\),\(\circled{5}\) 中选两个分别作为侧视图和俯视图,组成某个三棱锥的三视图,则所选侧视图和俯视图的编号依次为\fillinblank{}(写出符合要求的一组答案即可).

\FigureLayoutDeclare{img/2021/national_paper_B_science/q16_fig1.png}{12.0cm}{58bp 64bp 76bp 8bp}
\bitmapfigure[max width=8.2cm]{img/2021/national_paper_B_science/q16_fig1.png}
\end{problem}

\section{解答题}

\begin{problem}
某厂研制了一种生产高精产品的设备,为检验新设备生产产品的某项指标有无提高,用一台旧设备和一台新设备各生产了 \(10\) 件产品,得到各件产品该项指标数据如下:

\begin{center}
\begin{tabular}{|c|*{10}{c|}}
\hline
旧设备 & 9.8 & 10.3 & 10.0 & 10.2 & 9.9 & 9.8 & 10.0 & 10.1 & 10.2 & 9.7 \\
\hline
新设备 & 10.1 & 10.4 & 10.1 & 10.0 & 10.1 & 10.3 & 10.6 & 10.5 & 10.4 & 10.5 \\
\hline
\end{tabular}
\end{center}

旧设备和新设备生产产品的该项指标的样本平均数分别记为 \(x\) 和 \(y\),样本方差分别记为 \(s_1^2\) 和 \(s_2^2\).

\begin{enumerate}
\item 求 \(x\),\(y\),\(s_1^2\),\(s_2^2\);

\item 判断新设备生产产品的该项指标的均值较旧设备是否有显著提高(如果 \(\dfrac{y-x}{\sqrt{\left( s_1^2+s_2^2 \right)/10}}\ge 2\),则认为新设备生产产品的该项指标的均值较旧设备有显著提高,否则不认为有显著提高).
\end{enumerate}
\end{problem}

\begin{problem}
如图,四棱锥 \(P-ABCD\) 的底面是矩形,\(PD\perp\) 底面 \(ABCD\),\(PD=DC=1\),\(M\) 为 \(BC\) 的中点,且 \(PB\perp AM\).

\begin{enumerate}
\item 求 \(BC\);

\item 求二面角 \(A-PM-B\) 的正弦值.
\end{enumerate}

\bitmapfigure[max width=6.0cm]{img/2021/national_paper_B_science/q18_fig1.png}
\end{problem}

\begin{problem}
记 \(S_n\) 为数列 \(\{a_n\}\) 的前 \(n\) 项和,\(b_n\) 为数列 \(\{S_n\}\) 的前 \(n\) 项积,已知 \(\frac{2}{S_n}+\frac{1}{b_n}=2\).

\begin{enumerate}
\item 证明:数列 \(\{b_n\}\) 是等差数列;

\item 求 \(\{a_n\}\) 的通项公式.
\end{enumerate}
\end{problem}

\begin{problem}
设函数 \(f(x)=\ln(a-x)\),已知 \(x=0\) 是函数 \(y=xf(x)\) 的极值点.

\begin{enumerate}
\item 求 \(a\);

\item 设函数 \(g(x)=\dfrac{x+f(x)}{xf(x)}\),证明:\(g(x)<1\).
\end{enumerate}
\end{problem}

\begin{problem}
已知抛物线 \(C:x^2=2py\)(\(p>0\))的焦点为 \(F\),且 \(F\) 与圆 \(M:x^2+(y+4)^2=1\) 上点的距离的最小值为 \(4\).

\begin{enumerate}
\item 求 \(p\);

\item 若点 \(P\) 在 \(M\) 上,\(PA\),\(PB\) 是 \(C\) 的两条切线,\(A\),\(B\) 是切点,求 \(\triangle PAB\) 面积的最大值.
\end{enumerate}
\end{problem}

\section{选考题:解答题}

\begin{problem}
在直角坐标系 \(xOy\) 中,圆 \(C\) 的圆心为 \(C(2,1)\),半径为 \(1\).

\begin{enumerate}
\item 写出圆 \(C\) 的一个参数方程;

\item 过点 \(F(4,1)\) 作圆 \(C\) 的两条切线. 以坐标原点为极点,\(x\) 轴正半轴为极轴建立坐标系,求这两条切线的极坐标方程.
\end{enumerate}
\end{problem}

\begin{problem}
已知函数 \(f(x)=|x-a|+|x+3|\).

\begin{enumerate}
\item 当 \(a=1\) 时,求不等式 \(f(x)\ge 6\) 的解集;

\item 若 \(f(x)>-a\),求 \(a\) 的取值范围.
\end{enumerate}
\end{problem}

